

\section{Aplicaciones de Satisfactibilidad}
\subsection{Programación de exámenes universitarios}
Se deben programar tres exámenes en dos franjas horarias disponibles: \textbf{mañana} (\(t1\)) y \textbf{tarde} (\(t2\)). Los exámenes son:

\begin{itemize}
    \item \textbf{Matemáticas} (\(\boldsymbol{M}\))
    \item \textbf{Física} (\(\boldsymbol{P}\))
    \item \textbf{Ciencias de la Computación} (\(\boldsymbol{C}\))
\end{itemize}

Las restricciones son:
\begin{enumerate}
    \item Cada examen debe programarse en \textbf{exactamente una} franja horaria.
    \item Los exámenes de Matemáticas y Física \textbf{no pueden coincidir} en la misma franja.
    \item El examen de Ciencias de la Computación \textbf{debe programarse antes} que el de Matemáticas.
\end{enumerate}

\subsubsection*{Definición de Variables Booleanas}
Cada variable \( X_i \) indica si el examen \( X \) está programado en la franja \( ti \):

\begin{itemize}
    \item \( M_{1}, M_{2} \): Matemáticas en la franja $t1$ o $t2$.
    \item \( P_{1}, P_{2} \): Física en la franja $t1$ o $t2$.
    \item \( C_{1}, C_{2} \): Ciencias de la Computación en la franja $t1$ o $t2$.
\end{itemize}

\subsubsection*{Codificación en lógica proposicional}

\begin{enumerate}
\item Cada examen debe estar en una única franja
\begin{align*}
  (M_{1} \lor M_{2}) \land \neg (M_{1} \land M_{2}) \\
    (P_{1} \lor P_{2}) \land \neg (P_{1} \land P_{2}) \\
    (C_{1} \lor C_{2}) \land \neg (C_{1} \land C_{2})
\end{align*}

\item Matemáticas y Física no pueden coincidir
\begin{align*}
   \neg (M_{1} \land P_{1}), \quad \neg (M_{2} \land P_{2})\end{align*}

\item Ciencias de la Computación antes que Matemáticas 
\begin{align*}
    &C_{2} \rightarrow \neg M_{1}.\\
    &\text{(Si C está en $t2$, M no puede estar en $t1$)}
\end{align*}

\end{enumerate}


\subsubsection*{Solución SAT}
Una asignación válida que cumple todas las restricciones es:

\begin{align*}
    C_{1} &= \text{true}, \quad C_{2} = \text{false} \\
    M_{2} &= \text{true}, \quad M_{1} = \text{false} \\
    P_{2} &= \text{true}, \quad P_{1} = \text{false}
\end{align*}

\subsubsection*{Interpretación de la Solución:}
\begin{itemize}
    \item \textbf{Ciencias de la Computación} → Mañana (\(T1\))
    \item \textbf{Matemáticas} → Tarde (\(T2\))
    \item \textbf{Física} → Tarde (\(T2\))
\end{itemize}

La asignación respeta todas las restricciones, garantizando que \textbf{Ciencias de la Computación ocurre antes que Matemáticas} y que \textbf{Matemáticas y Física no coinciden}.





\subsection{Juego del Gato}
En este problema, modelamos el juego del Gato (\textit{Tic-Tac-Toe}) como un problema de \textbf{satisfactibilidad (SAT)} en lógica proposicional. El objetivo es verificar si el jugador \(X\) tiene una \textbf{jugada ganadora} en su próximo turno.

\subsubsection*{Definición del Problema}

Dado un tablero \(3 \times 3\), donde algunas casillas ya están ocupadas por \(X\) o \(O\), queremos determinar si existe una celda vacía donde \(X\) pueda jugar para ganar inmediatamente.

\subsubsection*{Representación del Tablero}

Definimos el estado de cada celda \((i,j)\) usando dos variables booleanas:

\begin{itemize}
    \item \( X_{i,j} \): La celda \((i,j)\) está ocupada por \(X\).
    \item \( O_{i,j} \): La celda \((i,j)\) está ocupada por \(O\).
\end{itemize}

Las celdas vacías no se modelan explícitamente, sino que se deducen por la ausencia de \(X_{i,j}\) y \(O_{i,j}\).

\subsubsection*{Restricciones del Problema}

\begin{enumerate}
\item Cada celda debe estar en exactamente un estado, lo que se expresa con la siguiente restricción:

\begin{equation}
    (\neg X_{i,j} \land \neg O_{i,j}) \lor \left( (X_{i,j} \lor O_{i,j})  \land 
    \neg(X_{i,j} \land O_{i,j}) \right), \quad \forall i,j.
\end{equation}

Esta ecuación garantiza que cada celda la celda esté vacía o esté ocupada por \(X\) o \(O\), pero no por ambos a la vez.

\item Condición 1: \(X\) debe jugar en una celda vacía

Para que \(X\) pueda realizar una jugada, debe colocar su marca en una celda que aún no esté ocupada. Dado que las celdas vacías se identifican como aquellas donde \(X_{i,j}\) y \(O_{i,j}\) son falsas, la condición de jugada se expresa como:

\begin{equation}
    \bigvee_{i,j} ((\neg X_{i,j} \land \neg O_{i,j}) \rightarrow X_{i,j}).
\end{equation}

Esta ecuación establece que si una celda no está ocupada por \(X\) ni \(O\), entonces \(X\) puede colocar su marca en ella.

\item Condición 2: \(X\) debe completar una fila, columna o diagonal.


\begin{itemize}
    \item \textbf{Fila completa:}
    \begin{equation}
        (X_{1,1} \land X_{1,2} \land X_{1,3}) \lor
        (X_{2,1} \land X_{2,2} \land X_{2,3}) \lor
        (X_{3,1} \land X_{3,2} \land X_{3,3}).
    \end{equation}

    \item \textbf{Columna completa:}
    \begin{equation}
        (X_{1,1} \land X_{2,1} \land X_{3,1}) \lor
        (X_{1,2} \land X_{2,2} \land X_{3,2}) \lor
        (X_{1,3} \land X_{2,3} \land X_{3,3}).
    \end{equation}

    \item \textbf{Diagonal completa:}
    \begin{equation}
        (X_{1,1} \land X_{2,2} \land X_{3,3}) \lor
        (X_{1,3} \land X_{2,2} \land X_{3,1}).
    \end{equation}
\end{itemize}

\end{enumerate}

\subsubsection*{Ejemplo de Tablero y Evaluación}

Supongamos que el tablero actual es:

\begin{center}
\begin{tabular}{|c|c|c|}
    \hline
    \(X\) & \(X\) & \(\emptyset\) \\ \hline
    \(O\) & \(O\) & \(X\) \\ \hline
    \(\emptyset\) & \(O\) & \(\emptyset\) \\ \hline
\end{tabular}
\end{center}

En este caso:
- \(X\) puede jugar en la celda \( (1,3) \) y completar la primera fila.
- Como la fórmula es satisfactible, se confirma que \(X\) tiene una jugada ganadora.





\section{Deducción Natural}



\begin{tcolorbox}[colframe=blue!60, colback=blue!5, title=\textbf{Métodos de Demostración en Lógica Proposicional}]
Cuando queremos demostrar que una afirmación es verdadera, podemos utilizar distintos enfoques. En lógica proposicional, los dos principales métodos son:
\begin{itemize}
    \item \textbf{Deducción Natural:} Un sistema de reglas que nos permite derivar proposiciones paso a paso a partir de premisas iniciales.
    \item \textbf{Tablas de Verdad:} Un método exhaustivo que evalúa todas las combinaciones posibles de valores de verdad para determinar si una proposición es válida en todos los casos.
\end{itemize}
\end{tcolorbox}

Ambos métodos tienen ventajas y limitaciones, y en esta sección nos centraremos en la \textbf{deducción natural}.

\subsection{¿Qué es la Deducción Natural?}

La \textbf{deducción natural} es un sistema formal que imita la forma en que razonamos de manera estructurada: partiendo de \textbf{premisas}, aplicamos \textbf{reglas de inferencia} para derivar conclusiones que serán correctas siempre que las premisas lo sean. Este proceso exige un enfoque \textbf{creativo y metódico}, parecido a la resolución de problemas en programación o demostración de teoremas, pues no siempre es evidente qué regla usar ni en qué orden. 

A diferencia de las \textbf{tablas de verdad}, que comprueban la validez de una proposición al evaluar todas las combinaciones posibles de valores de verdad, la deducción natural construye el razonamiento paso a paso, ofreciendo una comprensión detallada de cómo se encadenan las inferencias hasta llegar a la conclusión.

\subsubsection*{Ventajas y Desventajas de la Deducción Natural}
\begin{itemize}
    \item \textbf{Ventajas:}
    \begin{itemize}
        \item Permite demostrar teoremas de manera sistemática y organizada.
        \item Es aplicable a lógica de primer orden y sistemas más complejos.
        \item Es proceso generativo que permite descubrir nuevas conclusiones.
        \item Es más eficiente que las tablas de verdad para proposiciones con muchas variables.
    \end{itemize}

    \item \textbf{Desventajas:}
    \begin{itemize}
        \item Requiere práctica y habilidad para construir demostraciones.
        \item No proporciona una forma directa de comprobar si una proposición es una tautología.
    \end{itemize}
\end{itemize}

\bigskip

A continuación, exploraremos algunas reglas esenciales de la deducción natural, junto con ejemplos prácticos.


\subsection{Reglas Básicas de la Deducción Natural}

\subsubsection{Regla de la Conjunción (\(\wedge\))}

\noindent \textbf{Ejemplo:} Supongamos que sabemos que:
\begin{itemize}
    \item \textbf{Pedro tiene una bicicleta.}
    \item \textbf{Pedro tiene un casco.}
\end{itemize}
De esto, podemos concluir:
\begin{center}
    \textbf{Pedro tiene una bicicleta y un casco.}
\end{center}
Esto es un caso de la \textbf{introducción de la conjunción}.

\begin{tcolorbox}[colframe=blue!60, colback=blue!5, title=\textbf{Regla de Introducción de la Conjunción}]
Si tenemos dos afirmaciones \(\phi\) y \(\psi\), podemos combinarlas con \(\land\).
\[
\frac{\phi \quad \psi}{\phi \wedge \psi} (\wedge i)
\]
\end{tcolorbox}
Sobre las líneas están las premisas de la regla.  Abajo de la línea se encuentra la conclusión.  
A la derecha de la regla tenemos el nombre de la regla abreviado.

\noindent \textbf{Ejemplo:}  
Si sabemos que:
\[
\text{Ana es ingeniera y Ana es programadora.}
\]
Podemos concluir:
\[
\text{Ana es programadora.}
\]
\begin{tcolorbox}[colframe=red!60, colback=red!5, title=\textbf{Regla de Eliminación de la Conjunción}]
Si sabemos que algo es cierto en conjunto (\(\phi \wedge \psi\)), podemos extraer cualquiera de sus partes.
\[
\frac{\phi \wedge \psi}{\phi} (\wedge e_1) \quad
\frac{\phi \wedge \psi}{\psi} (\wedge e_2)
\]
\end{tcolorbox}

\noindent \textbf{Ejercicio:} 
Queremos probar:
\[
p \wedge q, r \vdash q \wedge r
\]

Representación gráfica de la demostración:
\[
\Rulee{\Rulee{p \wedge q}{q}{\wedge e_2} \qquad r}{
q \wedge r}{\wedge i}
\]


\subsubsection{Regla de la Doble Negación}

\textbf{Motivación:}  
Si alguien dice:  
\begin{center}
    "\textbf{No es falso que Pedro es honesto.}"
\end{center}
¿Significa esto que Pedro es honesto?  
Sí. Esto es un caso de eliminación de la doble negación.

\begin{tcolorbox}[colframe=blue!60, colback=blue!5, title=\textbf{Regla de la Doble Negación}]
Una afirmación es equivalente a su doble negación.
\[
\frac{\neg \neg \phi}{\phi} (\neg \neg e) \quad
\frac{\phi}{\neg \neg \phi} (\neg \neg i)
\]
\end{tcolorbox}

\subsubsection{Regla de Eliminación de la Implicación (Modus Ponens)}

\noindent \textbf{Ejemplo:}  
Sabemos que:
\begin{itemize}
    \item Si estudio, aprobaré el examen.
    \item Estudié.
\end{itemize}
Podemos concluir:
\begin{center}
    \textbf{Aprobaré el examen.}
\end{center}

Esto es un caso del \textbf{Modus Ponens}.


\begin{tcolorbox}[colframe=blue!60, colback=blue!5, title=\textbf{Modus Ponens (\(\rightarrow e\))}]
Si tenemos \(\phi\) y \(\phi \rightarrow \psi\), podemos deducir \(\psi\).
\[
\frac{\phi \quad \phi \rightarrow \psi}{\psi} (\rightarrow e)
\]
\end{tcolorbox}

\begin{tcolorbox}[colframe=red!60, colback=red!5, title=\textbf{Modus Tollens (MT)}]
Si tenemos \(\phi \rightarrow \psi\) y \(\neg \psi\), podemos deducir \(\neg \phi\).
\[
\frac{\phi \rightarrow \psi \quad \neg \psi}{\neg \phi} (MT)
\]
\end{tcolorbox}



\subsubsection{Regla de Introducción de la Implicación}

\textbf{Ejemplo:}  
Supongamos que queremos demostrar la siguiente afirmación:  
\begin{center}
    \textbf{Si Juan estudia, aprobará el examen.}
\end{center}
¿Cómo lo haríamos?  
Podemos asumir temporalmente que Juan estudia, y si logramos demostrar que en ese caso aprueba el examen, entonces la afirmación original es válida.

\begin{tcolorbox}[colframe=blue!60, colback=blue!5, title=\textbf{Regla de Introducción de la Implicación}]
Para probar \(\phi \rightarrow \psi\), asumimos temporalmente \(\phi\) y demostramos \(\psi\).
\[
\frac{\ovalbox{\begin{tabular}{c}$\phi$\\ $\vdots$\\ $\psi$ \end{tabular}
               }}{\phi \rightarrow \psi}{(\rightarrow i)}
\]
\end{tcolorbox}
La regla expresa que para probar $\phi \rightarrow \psi$ debemos asumir temporalmente $\phi$ y  entonces probar $\psi$.  
En la prueba de $\psi$ se pueden usar cualquiera de las fórmulas obtenidas hasta el momento,  como premisas y conclusiones provisionales.

\noindent \textbf{Ejercicio:} 
Queremos demostrar \[
p \rightarrow q \vdash \neg q \rightarrow \neg p
\]
La mecánica de la regla es más complicada que las anteriores.
\bigskip

\begin{tabular}{p{4mm}p{40mm}p{25mm}}
1    & $p \rightarrow q$                 & premisa \\
\framebox{\begin{tabular}{@{}p{4mm}p{40mm}p{30mm}}
  2   & $\neg q$       & supuesto \\
  3   & $\neg p$       & MT 1,2\\
          \end{tabular}
         } \\[5mm]
4     & $\neg q \rightarrow \neg p$  & $\rightarrow i$ 2--3
\end{tabular}\\[4mm]

El cuadro sirve para delimitar el alcance del supuesto temporal  $\neg q$.
No podemos usar suposiciones dentro del cuadro en reglas fuera del cuadro.


\subsubsection{Reglas de la Disyunción (\(\vee\))}

\begin{tcolorbox}[colframe=blue!60, colback=blue!5, title=\textbf{Introducción de la Disyunción}]
Podemos introducir una disyunción a partir de cualquier proposición:
\[
\frac{\phi}{\phi \vee \psi} (\vee i_1) \quad
\frac{\psi}{\phi \vee \psi} (\vee i_2)
\]
\end{tcolorbox}

\bigskip

Supongamos que queremos probar una proposición $\chi$ asumiendo que $\phi \vee \psi$.
Puesto que no sabemos que argumento de la disyunción es verdadero, debemos obtener dos pruebas por separado que deberemos combinar en un solo argumento:
\begin{tcolorbox}[colframe=red!60, colback=red!5, title=\textbf{Eliminación de la Disyunción}]
Si queremos probar \(\chi\) asumiendo \(\phi \vee \psi\), necesitamos dos pruebas separadas para cada caso:
\[
\frac{\phi \vee \psi \quad \ovalbox{\begin{tabular}{c}$\phi$\\ $\vdots$\\ $\chi$ 
\end{tabular}
} 
\quad \ovalbox{\begin{tabular}{c}$\psi$\\ $\vdots$\\ $\chi$ 
\end{tabular}
}
}{\chi}{(\vee e)}
\]
\end{tcolorbox}



\subsubsection{Regla de la Negación}

Una \textbf{contradicción} en lógica proposicional se expresa como \(\phi \land \neg \phi\), donde \(\phi\) es cualquier fórmula. A esta contradicción se le denomina \(\bot\). La regla de la negación describe cómo usar una contradicción y cómo introducir una negación al demostrar que un supuesto conduce a \(\bot\).

\begin{tcolorbox}[colframe=blue!60, colback=blue!5, title=\textbf{Regla de la Negación}]
\[
\Rulee{\phi \quad \neg \phi}{\bot}{\neg e}
\quad
\Rulee{\bot}{\phi}{\bot e}
\quad
\Rulee{
\boxed{
\begin{matrix}
\phi \\
\vdots \\
\bot
\end{matrix}
}
}{\neg \phi}{\neg i}
\]
\end{tcolorbox}

\begin{itemize}
    \item \(\neg e\) (eliminación de la negación) nos dice que si \(\phi\) y \(\neg \phi\) se derivan simultáneamente, obtenemos \(\bot\). 
    \item \(\bot e\) nos indica que, a partir de \(\bot\), podemos concluir \textit{cualquier} fórmula (principio de explosión).
    \item \(\neg i\) (introducción de la negación) expresa que, si al suponer \(\phi\) derivamos una contradicción (\(\bot\)), entonces concluimos \(\neg \phi\). 
\end{itemize}



\bigskip

\subsubsection*{Ejercicio Final: Demostración con Disyunciones y Conjunciones}

Veamos un ejemplo que combina varias reglas de inferencia para demostrar:

\[
p \wedge (q \vee r) \;\;\vdash\;\; (p \wedge q) \;\vee\; (p \wedge r).
\]

La idea general es extraer \(p\) y la disyunción \(q \vee r\) de la premisa, y luego considerar \textbf{dos casos} (uno donde \(q\) es verdadero, y otro donde \(r\) es verdadero) para demostrar la disyunción final. La prueba se detalla a continuación:

% \bigskip

\begin{tabular}{p{4mm}p{50mm}p{25mm}}
1 & $p \wedge ( q \vee r)$ & \textit{premisa} \\ 
2 & $p$                    & $\wedge e_1$ (de 1)  \\
3 & $q \vee r$             & $\wedge e_2$ (de 1)  \\
\framebox{\begin{tabular}{@{}p{4mm}p{46mm}p{30mm}}
4  & $q$              & \textit{supuesto}\\
5  & $p \wedge q$     & $\wedge i$ (de 2,4)\\
6  & $(p \wedge q) \vee (p \wedge r)$ & $\vee i_1$ (de 5)
\end{tabular}
} \\[5mm]
\framebox{\begin{tabular}{@{}p{4mm}p{46mm}p{30mm}}
7  & $r$             & \textit{supuesto}\\
8  & $p \wedge r$    & $\wedge i$ (de 2,7)\\
9  & $(p \wedge q) \vee (p \wedge r)$ & $\vee i_2$ (de 8)
\end{tabular}
} \\[5mm]
10 & $(p \wedge q) \vee (p \wedge r)$ & $\vee e$ (3, 4--6, 7--9)
\end{tabular}

\subsection*{Conclusión}

La deducción natural es más que un conjunto de reglas: se trata de un sistema de prueba que garantiza, por un lado, la \textbf{solidez} (no se puede demostrar algo falso) y, por otro, la \textbf{completitud} (todo lo que es válido puede probarse dentro del sistema), al menos en el contexto de la lógica proposicional. 
Estas dos propiedades son cruciales para confiar en cualquier sistema lógico: nos aseguran que no probaremos fórmulas falsas y que no dejaremos sin demostrar las verdaderas.


\begin{tcolorbox}[colframe=orange!60, colback=orange!5,  title=\textbf{Preguntas para Reflexionar}]
\begin{itemize}
    \item \textbf{Solidez (Soundness):}  
    Si podemos derivar \(\alpha\) a partir de \(\Gamma\) mediante deducción natural, entonces \(\alpha\) es una consecuencia lógica de \(\Gamma\) en todas las interpretaciones.  
    En términos formales:  
    \[
    \Gamma \vdash \alpha 
    \quad \Longrightarrow \quad
    \Gamma \models \alpha.
    \]
   

    \item \textbf{Completitud (Completeness):}  
    Si \(\alpha\) es una consecuencia lógica de \(\Gamma\) (es decir, verdadera en todas las interpretaciones donde \(\Gamma\) lo es), entonces también podemos derivar \(\alpha\) dentro del sistema de deducción natural.
    En términos formales:  
    \[
    \Gamma \models \alpha 
    \quad \Longrightarrow \quad
    \Gamma \vdash \alpha.
    \]
\end{itemize}
\end{tcolorbox}
% Gracias a estas propiedades, la deducción natural logra un balance entre \emph{confiabilidad} y \emph{alcance}: su \textbf{consistencia} garantiza que no introducimos falsedades en nuestras derivaciones, mientras que su \textbf{completitud} nos asegura que, de existir un camino lógico para demostrar una proposición válida, podremos hallarlo dentro de este sistema.

Este \textbf{enfoque riguroso y ordenado} para el razonamiento lógico tiene implicaciones tanto en el ámbito académico—donde se emplea para demostrar teoremas y validar argumentos formales—como en la vida diaria, al entrenar nuestras habilidades de \emph{pensamiento analítico}. Aun si no aplicamos conscientemente el lenguaje lógico en cada situación, los principios de consistencia y completitud ofrecen un horizonte de \emph{certeza} para saber que los razonamientos formales se construyen de forma sólida y exhaustiva.

% \begin{tcolorbox}[colframe=orange!60, colback=orange!5, title=\textbf{Preguntas para Reflexionar}]
% \begin{itemize}
%     \item ¿Cómo aplicarías estas reglas a problemas de la vida real, donde las premisas no siempre están tan claras?
%     \item ¿En qué ocasiones crees que la verificación exhaustiva (al estilo de las tablas de verdad) podría resultar más práctica que la deducción paso a paso?
%     \item ¿De qué manera los conceptos de \emph{consistencia} y \emph{completitud} refuerzan la confianza en tus razonamientos formales?
% \end{itemize}
% \end{tcolorbox}




